\documentclass[11pt]{article}
\usepackage[a4paper,margin=1in]{geometry}
\usepackage{times}
\usepackage[utf8]{inputenc}
\usepackage[T1]{fontenc}
\usepackage{hyperref}
\usepackage{enumitem}
\usepackage{titlesec}
\usepackage{tcolorbox}
\usepackage{fancyhdr}
\usepackage{parskip}
\usepackage{lastpage}

\hypersetup{colorlinks=true,linkcolor=blue,urlcolor=blue,citecolor=blue}
\titleformat{\section}{\large\bfseries}{\thesection}{1em}{}
\pagestyle{fancy}
\fancyhf{}
\cfoot{\thepage\ / \pageref{LastPage}}

\title{\textbf{Playing Through an AI Incident}}
\author{Sofia G. Gallego \\ \textit{Cosmoventures, Paris, France}}
\date{}

\begin{document}
\maketitle
\thispagestyle{fancy}

\begin{tcolorbox}[colback=gray!5,colframe=gray!50,title=\textbf{Abstract}]
\emph{Little Agent Lab} is a free browser game that teaches three ideas central to AI oversight --- capability separation, gated review, and the fact that a labeled safeguard may not actually check anything --- through three short, escalating missions built around a single, always-visible goal. Every ``agent'' is a small deterministic state machine, not a live model, so behavior is fully reproducible and safe for a non-expert, including a child, to inspect directly rather than read about. The purpose is early awareness: giving the next generation of users, well before they encounter these ideas in a policy document or a headline, a concrete, hands-on sense of what ``oversight'' and ``verification'' actually mean. This is a small, work-in-progress idea we plan to keep extending. We describe what it does, why we built it that way, and what we'd add next.
\end{tcolorbox}

\section{What it is}
\emph{Little Agent Lab} is a browser game with no accounts, no network calls, and no live model --- everything is a small, deterministic, hand-scripted simulation. You pick a builder agent, add up to two helper agents, set a gate that controls outside access, and run the plan. What actually happens plays out as a short, literal trace of events, so you can watch --- not just be told --- what each agent did.

Three missions, unlocked in sequence:
\begin{itemize}[topsep=2pt,itemsep=2pt]
\item \textbf{Weather Lantern.} Two valid strategies exist: give the builder a helper whose route structurally never touches the private item, or open the gate but pair it with an inspector that actually checks what leaves. First taste of capability separation.
\item \textbf{Garden Delivery.} No safe shortcut this time --- the cautious agent simply can't finish the job alone. The only way to win is to take the riskier route and pair it with a real check.
\item \textbf{Message Relay.} Two helpers offer to guard the gate. Only one of them actually blocks anything, and you can't tell which from its description --- you have to run the plan and read the trace to find out. This is the core idea of the whole project: a safeguard is only as good as what it does, not what it's called.
\end{itemize}

\section{Why}
Most explanations of AI oversight are prose aimed at people who already work in the field. This project is aimed the other way: at a general audience, and especially at kids, who will grow up as the users of these systems rather than their builders. The bet is that a five-minute playable moment --- picking the wrong-looking guard and watching it fail to stop anything --- builds an intuition that a paragraph doesn't. That intuition, at scale, is a small but real input into how the next generation of people reasons about AI systems they didn't build and don't fully see inside.

\section{Status and what's next}
This is a solo, short-timeframe build, iterated a few times based on direct feedback rather than a formal study: an earlier, more abstract version of the same idea (an unlabeled stream of requests to allow or block) was tried and dropped because it needed its own tutorial and stopped feeling like a game. The current version is simpler on purpose. All three missions' logic is covered by automated tests (27, all passing), and every screen has been checked in a real browser for rendering and console errors --- but no real player, child or adult, has used it yet, and that's the next step. Planned extensions: more missions, a version of the mechanic where the fake safeguard's tell changes each time so memorizing the answer stops working, and an actual playtest with the age group it's meant for.

\begin{tcolorbox}[colback=orange!5,colframe=orange!50!black,title=\textbf{Limitations and Dual-Use Considerations}]
\textbf{Limitations.} No human-subject testing has been done; every usability decision so far came from one non-blinded evaluator. Three missions is enough to show the mechanic, not enough to know if the lesson generalizes. The game's five-line trace is far more legible than a real system's actual safeguards, and a player could walk away overconfident about how easy real verification is. The implementation was AI-assisted (Claude Code), disclosed here per the sprint's guidance on describing contributions accurately.

\textbf{Dual-use.} Risk is low. Every agent name, destination, and payload in the game is fictional; nothing represents a real API, protocol, or system. The one mechanic that resembles a ``bypass'' (Message Relay's fake guard) is a fixed, pre-authored branch chosen from two labeled cards, not a generative or discoverable exploit --- there is nothing here a player could apply against a real system. No model was jailbroken and no organization's real systems were touched in building this.
\end{tcolorbox}

\vspace{0.5em}
\footnotesize
Project source: \url{https://github.com/nefinia/little-agent-lab} (public, playable via the repo's GitHub Pages workflow).

\end{document}
